%%%%%%%%%%%%%%%%%%%%%%%%%%%%%%%%%%%%%%%%%%%%%%%%%%%%%%%%%%
% LMTH 2040 — Calculus: A Multidisciplinary Approach
% Homework 01
% Fall 2026
%%%%%%%%%%%%%%%%%%%%%%%%%%%%%%%%%%%%%%%%%%%%%%%%%%%%%%%%%%

\documentclass[11pt]{article}

% Adjust path as needed
\usepackage{../langcalc}

\begin{document}

\courseheader
    {Homework 01: Patterns and Accumulation}
    {LMTH 2040 --- Calculus: A Multidisciplinary Approach}
    {Fall 2026}
    {Due Monday, August 31}

\begin{keyidea}
This assignment asks you to investigate patterns before being given
formulas for them. Draw pictures, calculate examples, experiment with
technology, make conjectures, and explain why you think your conclusions
are true.
\end{keyidea}

\section*{Problem 1 \quad Triangular Numbers}

Consider the following arrangements of dots:

\begin{center}
\begin{tabular}{cccc}
Stage 1 & Stage 2 & Stage 3 & Stage 4 \\[6pt]

$\bullet$
&
\begin{tabular}{c}
$\bullet$\\[-2pt]
$\bullet\ \bullet$
\end{tabular}
&
\begin{tabular}{c}
$\bullet$\\[-2pt]
$\bullet\ \bullet$\\[-2pt]
$\bullet\ \bullet\ \bullet$
\end{tabular}
&
\begin{tabular}{c}
$\bullet$\\[-2pt]
$\bullet\ \bullet$\\[-2pt]
$\bullet\ \bullet\ \bullet$\\[-2pt]
$\bullet\ \bullet\ \bullet\ \bullet$
\end{tabular}
\end{tabular}
\end{center}

The number of dots in the \(n\)th arrangement is called the
\(n\)th \emph{triangular number}, which we will denote by \(T_n\).

\begin{enumerate}[label=\textbf{\alph*.}, leftmargin=2em]

\item
Draw the next three stages and determine \(T_5\), \(T_6\), and \(T_7\).

\item
Find a way to arrange, duplicate, reflect, cut, or otherwise transform
one of these pictures so that you can determine \(T_n\) without adding

\[
1+2+3+\cdots+n
\]

one term at a time.

Include a picture of your construction.

\item
Use your construction to find a formula for \(T_n\).

Explain why your formula works for \emph{every} positive integer \(n\),
not only for the examples you have drawn.

\item
Now consider the ratio

\[
\frac{T_n}{n^2}.
\]

Calculate this ratio for several values of \(n\).

What seems to happen as \(n\) becomes very large?

Make a conjecture and then use your formula from part (c) to explain
whether your conjecture must be true.

\end{enumerate}

\begin{reflection}
\textbf{Computational Experiment.}

Use a computational tool of your choice to extend your investigation
beyond what would be convenient to calculate by hand.

You might use Python or Jupyter, a spreadsheet, Desmos, GeoGebra,
Mathematica, a calculator, or another computational tool.

Investigate

\[
\frac{T_n}{n^2}
\]

for increasingly large values of \(n\). You might construct a table,
make a graph, write a short program, or develop another useful
representation.

Include enough information that someone else could reproduce your
experiment.

Briefly explain:

\begin{enumerate}[label=\arabic*., leftmargin=2em]
    \item what tool you used and what you asked it to do;
    \item what pattern you observed; and
    \item what the computation shows---and what it does \emph{not}
    establish by itself.
\end{enumerate}
\end{reflection}


% =========================================================
\section*{Problem 2 \quad A Growing Shape}

Consider the following pattern made from unit squares.

\begin{center}
\renewcommand{\arraystretch}{0.8}

\begin{tabular}{ccc}

\textbf{Stage 1}
&
\textbf{Stage 2}
&
\textbf{Stage 3}
\\[10pt]

\(\blacksquare\)

&

\begin{tabular}{c}
\(\phantom{\blacksquare}\blacksquare\phantom{\blacksquare}\)\\
\(\blacksquare\blacksquare\blacksquare\)\\
\(\phantom{\blacksquare}\blacksquare\phantom{\blacksquare}\)
\end{tabular}

&

\begin{tabular}{c}
\(\phantom{\blacksquare\blacksquare}\blacksquare
  \phantom{\blacksquare\blacksquare}\)\\
\(\phantom{\blacksquare}
  \blacksquare\blacksquare\blacksquare
  \phantom{\blacksquare}\)\\
\(\blacksquare\blacksquare\blacksquare
  \blacksquare\blacksquare\)\\
\(\phantom{\blacksquare}
  \blacksquare\blacksquare\blacksquare
  \phantom{\blacksquare}\)\\
\(\phantom{\blacksquare\blacksquare}\blacksquare
  \phantom{\blacksquare\blacksquare}\)
\end{tabular}

\end{tabular}
\end{center}

Let \(S_n\) denote the number of unit squares in Stage \(n\).

\begin{enumerate}[label=\textbf{\alph*.}, leftmargin=2em]

\item
Draw Stages 4 and 5.

Determine

\[
S_1,\quad S_2,\quad S_3,\quad S_4,\quad S_5.
\]

\item
Instead of looking only at the total number of squares, investigate
how the figure \emph{grows}.

How many new squares are added in passing from one stage to the next?

Describe the pattern you find.

\item
Find a formula for \(S_n\).

Your explanation should make use of the geometry of the figure rather
than simply presenting an algebraic formula.

\item
Compare the sequences \(T_n\) and \(S_n\).

What is similar about the way they grow?

What is different?

Can one of the sequences help you understand the other?

\end{enumerate}

\begin{reflection}
\textbf{Computational Experiment.}

Use a computational tool of your choice to investigate the pattern at
scales that would be inconvenient to draw by hand.

Your goal is \emph{not} simply to generate more terms. Use computation
to discover something about the pattern.

For example, you might investigate:

\begin{itemize}[leftmargin=2em]
    \item the total number of squares \(S_n\);
    \item the number of new squares added at each stage;
    \item ratios involving \(S_n\);
    \item comparisons between \(S_n\), \(T_n\), \(n\), and \(n^2\);
    \item how quickly different measurements grow;
    \item or another pattern that you notice.
\end{itemize}

Present one computational observation that you find mathematically
interesting.

State a conjecture suggested by your experiment. Then explain what
additional mathematical reasoning would be necessary to establish
that conjecture.
\end{reflection}


% =========================================================
% =========================================================
\section*{Reading \quad Otto Toeplitz}

Read the \textbf{Preface to the English Edition} and the opening
section of Part I at least through the second paragraph on page 6 of \textbf{``The Nature of the Infinite Process,''}
from Otto Toeplitz, \emph{The Calculus: A Genetic Approach}.

\begin{historicalnote}
Toeplitz describes this book as an attempt to develop a
\emph{genetic approach} to calculus. As you read, pay attention to
what he thinks is unsatisfactory about the usual presentation of
mathematics and what he proposes doing differently.
\end{historicalnote}

\subsection*{Response}

Write approximately one paragraph responding to the following:

\begin{quote}
What does Toeplitz mean by a ``genetic approach'' to calculus?

Identify one claim he makes about how mathematics should be learned
or presented. Explain the claim in your own words and then respond
to it: what do you find convincing, questionable, or unresolved?

Where useful, draw on your experience with one of the mathematical
problems in this assignment.
\end{quote}


% =========================================================
\section*{Submission}

Submit a single PDF containing your mathematical work, figures,
computational experiment, and reading response.

Your solutions should communicate your reasoning, not only your final
answers. Include diagrams when they are part of your argument.

If you use computational tools, briefly
identify the tool and describe how you used it. You are responsible
for understanding and explaining any mathematics, code, or claims
included in your submission.

\end{document}